\section{Discussion}
% OUTLINE -- draft next. No prose yet.
% - What the horizon guarantee buys: the per-cycle certificate \citep{busico_m4} becomes a
%   runtime, horizon-wide safety monitor -- a long-run frequency guarantee on the joint event.
% - The certificate's operational value, quantified (Fig. 1): meeting a quality spec is not
%   safety; the health-blind loop is "in spec" yet structurally unsafe over the horizon. The
%   psi-budget is what converts a per-cycle bound into a sustained one.
% - Deadtime (Fig. 2 + O1): the cost of delayed feedback is the coverage RATE, not validity; the
%   selection penalty stays zero.
% - Adaptive vs static calibration (Fig. 3): a fixed union bound cannot buy horizon coverage
%   without going vacuous; adaptivity is necessary at realistic K.
% - Limitations: reduced-order (Hammerstein) dynamics, not a tray-resolved transient; illustrative
%   time constants to be pinned per column; DWSIM-dynamic high-fidelity confirmation is future work
%   (the in-sandbox loop proves the distribution-free claim, which is fidelity-agnostic for validity);
%   the deadtime rate-transient is small in this well-behaved loop, and a drift stress test would
%   make it visible.
% - Bridge to plantwide: the eps-budget allocated across units (modular risk allocation) is the
%   next step that scales the certificate from one column to a plant.
